Fast phylogeny reconstruction has changed since our chapter on fast
phylogeny reconstruction was published five years
ago~\cite{hau21:fas}. The changes concern not only the technical side
of distance computation between sequences, but also how to best access
genome sequences. Our forthcoming revised chapter reflects both the
changes in distance computation and in data acquisition.

Our new chapter contains a tutorial centered on the genomes of a
commensal serovar of \emph{Escherichia coli}, O16:H48, which is
closely related to the model strain K12. We reconstruct the phylogeny
of its genomes, and place the phylogeny in the context of almost 7,000
complete \emph{E. coli} genomes. The \ty{free} repo contains code for
setting up and following the tutorial. For set up, follow the
instructions in the README. In the following sections we describe the
three auxiliary scripts used during the tutorial.
